\documentclass[../../main.tex]{subfiles}% 注意这里的文件路径不能用 ./main.tex ,否则用latexmk编译子文件会报错
\graphicspath{{\subfix{./image/}}{\subfix{../../image/}}} % 指定图片引用路径,后续可以直接使用图片文件名
% 如果图片存放在上级目录,则使用 ../../image/ 作为备用路径

% 例如:
% \begin{figure}[H]
% \centering
% \includegraphics[width=0.3\textwidth]{图.png}
% \caption{}
% \label{figure:图}
% \end{figure}
% 注意:上述\label{}一定要放在\caption{}之后,否则引用图片序号会只会显示??.

\begin{document}

\section{欧拉函数}

\begin{definition}[欧拉函数]
\textbf{欧拉函数} $\varphi(m)$ 定义为：对于正整数 $m$，$\varphi(m)$ 表示不超过 $m$ 且与 $m$ 互素的正整数的个数。特别地，规定 $\varphi(1)=1$.
\end{definition}

\begin{theorem}\label{theorem:欧拉函数的基本性质}
设$\varphi$为欧拉函数,$p,m,n$为正整数,则
\begin{enumerate}[(1)]
\item 若 $p$ 为素数，则 $\varphi(p)=p-1$，且 $\varphi(p^\alpha)=p^\alpha-p^{\alpha-1}$（$\alpha\geqslant 1$）.

\item 若 $m,n$ 互素，则 $\varphi(mn)=\varphi(m)\varphi(n)$.

\item 若 $n=p_1^{\alpha_1}\cdots p_r^{\alpha_r}$ 为 $n$ 的标准分解式，则
\begin{align*}
\varphi(n)=n\prod_{i=1}^r\left(1-\frac1{p_i}\right).
\end{align*}

\item\label{theorem:欧拉函数的基本性质-4} $\mathbf{Gauss}$\textbf{恒等式}:
\begin{align}\label{eq:euler_sum}
\sum_{d\mid n}\varphi(d)=n.
\end{align}
\end{enumerate}
\end{theorem}
\begin{proof}
\begin{enumerate}[(1)]
\item 当 $p$ 为素数时，$1$ 到 $p$ 中只有 $p$ 与 $p$ 不互素，故 $\varphi(p)=p-1$. 对 $p^\alpha$，在 $1$ 到 $p^\alpha$ 中与 $p^\alpha$ 不互素的数恰好是 $p$ 的倍数，共有 $p^{\alpha-1}$ 个，因此 $\varphi(p^\alpha)=p^\alpha-p^{\alpha-1}$.

\item 设 $m,n$ 互素. 将 $1$ 到 $mn$ 的数按模 $m$ 和模 $n$ 的剩余类对应，由中国剩余定理，$k$ 与 $mn$ 互素当且仅当 $k$ 与 $m$ 互素且 $k$ 与 $n$ 互素，因此 $\varphi(mn)=\varphi(m)\varphi(n)$.

\item 由 (1)(2)直接得到
\begin{align*}
\varphi(n)=\prod_{i=1}^r\varphi(p_i^{\alpha_i})
=\prod_{i=1}^r p_i^{\alpha_i}\left(1-\frac1{p_i}\right)
=n\prod_{i=1}^r\left(1-\frac1{p_i}\right).
\end{align*}

\item 考虑集合 $S=\{1,2,\dots,n\}$，显然 $|S|=n$. 对任意 $k\in S$，令 $d=\gcd(k,n)$，则 $d$ 是 $n$ 的正因数. 对每个 $d\mid n$，记
\begin{align*}
A_d=\{k\in S: \gcd(k,n)=d\}.
\end{align*}
则所有 $A_d$ 构成 $S$ 的一个划分，故
\begin{align*}
\sum_{d\mid n}|A_d|=n.
\end{align*}
下面计算 $|A_d|$. 若 $k\in A_d$，则 $k=dq$，$n=dm$，其中 $\gcd(q,m)=1$，且 $1\leqslant q\leqslant m$. 因此 $q$ 的个数恰好为 $\varphi(m)=\varphi(\frac{n}{d})$. 所以 $|A_d|=\varphi(\frac{n}{d})$. 于是
\begin{align*}
\sum_{d\mid n}\varphi\left(\frac nd\right)=n.
\end{align*}
当 $d$ 遍历 $n$ 的所有正因数时，$\frac{n}{d}$ 也遍历 $n$ 的所有正因数，故
\begin{align*}
\sum_{d\mid n}\varphi(d)=n.
\end{align*}
\end{enumerate}
\end{proof}








\end{document}